\documentclass[12pt, a4paper]{article}

\usepackage{geometry}
\geometry{margin=1.15in, top=1.2in, bottom=1.2in}

\usepackage{fontspec}
\setmainfont{TeX Gyre Pagella}
\usepackage{microtype}

\usepackage{amsmath, amsthm, amssymb, mathtools}

\usepackage{hyperref}
\hypersetup{colorlinks=true, linkcolor=black, citecolor=black, urlcolor=black}

\usepackage{booktabs, enumitem, xcolor, setspace}
\usepackage{abstract}

% ── Theorem environments ─────────────────────────────────────────────────────
\theoremstyle{plain}
\newtheorem{theorem}{Theorem}[section]
\newtheorem{proposition}[theorem]{Proposition}
\newtheorem{lemma}[theorem]{Lemma}
\newtheorem{corollary}[theorem]{Corollary}

\theoremstyle{definition}
\newtheorem{definition}[theorem]{Definition}
\newtheorem{example}[theorem]{Example}

\theoremstyle{remark}
\newtheorem{remark}[theorem]{Remark}

% ── Notation ─────────────────────────────────────────────────────────────────
\newcommand{\calX}{\mathcal{X}}
\newcommand{\calM}{\mathcal{M}}
\newcommand{\calF}{\mathcal{F}}
\newcommand{\calC}{\mathcal{C}}
\newcommand{\calB}{\mathcal{B}}
\newcommand{\calO}{\mathcal{O}}
\newcommand{\calR}{\mathcal{R}}
\newcommand{\bR}{\mathbb{R}}
\newcommand{\vpi}{\pi}
\newcommand{\ve}{\varepsilon}
\newcommand{\dM}{d_{\calM}}
\newcommand{\norm}[1]{\left\lVert #1 \right\rVert}
\newcommand{\inner}[2]{\langle #1,\, #2 \rangle}
\newcommand{\fiber}[1]{\vpi^{-1}(#1)}
\newcommand{\connD}{\nabla}
\newcommand{\curv}{\Omega}

\renewcommand{\abstractnamefont}{\normalfont\small\bfseries}
\renewcommand{\abstracttextfont}{\normalfont\small}

% ── Title ────────────────────────────────────────────────────────────────────
\title{\textbf{Document Perturbation as Constrained Navigation\\
on a Semantic Manifold}}
\author{Flyxion\\
\small Independent Researcher}
\date{}

% ════════════════════════════════════════════════════════════════════════════
\begin{document}
\maketitle
\thispagestyle{empty}

% ── Abstract ─────────────────────────────────────────────────────────────────
\begin{abstract}
\noindent
We introduce a framework for generating controlled document variations through
constrained navigation in semantic space.  Traditional document augmentation
relies on synonym replacement, paraphrasing, or stochastic language-model
sampling, providing little formal control over semantic drift.  We propose a
\emph{semantic-budget framework} in which document transformations are
perturbation operators acting on a semantic manifold $\calM$.  A document is
a point in an embedding space; the text file is merely one coordinate
representation of an underlying geometric object.  Admissible transformations
are constrained by a finite semantic distance budget, local curvature bounds,
and conservation laws preserving core concepts.  The framework is developed as
exploration of the fibers of a projection $\vpi : \calX \to \calM$, where
$\calX$ is the space of syntactically valid textual realizations and $\calM$
encodes semantic identity.  We introduce a semantic connection $\connD$ on the
bundle that governs admissible transport, a curvature tensor $\curv_P$ encoding
the rigidity of different discourse communities, and a semantic accessibility
entropy measuring the logarithmic volume of the reachable set.  Motivated by
colored-noise diffusion sampling \cite{davidson2026colorednoisediffusionsampling},
we prove that budget allocation weighted by a semantic completion field
$\gamma(c,t)$ maximizes semantic diversity subject to a finite perturbation
budget---constituting semantic colored noise.  Mutation profiles are modelled as
profile-dependent metric tensors $M_P$ and connection structures $\connD_P$,
enabling academic, legal, journalistic, technical, and literary registers to
define semantic equivalence according to genuinely different geometric structures.
\end{abstract}

\newpage
\tableofcontents
\newpage

% ════════════════════════════════════════════════════════════════════════════
\section{Introduction}
\label{sec:intro}
% ════════════════════════════════════════════════════════════════════════════

Controlled variation of natural-language documents underlies data augmentation
for low-resource learning, robustness evaluation for retrieval and
question-answering systems, deduplication benchmark construction, plagiarism
detection calibration, and synthetic parallel corpus generation.  Despite the
centrality of these applications, the dominant approach remains conceptually
simple: select words at random, substitute synonyms from a lexical database,
and assume the resulting text remains meaningfully close to the original.

The limitations of this approach are well understood.  WordNet-based substitution
\cite{miller1995wordnet} treats all synonyms as equiprobable and ignores
contextual fit, producing phrases such as \emph{the investigator cautiously
promulgated the constitution of the antique edifice} as a paraphrase of
\emph{the scientist carefully published the structure of the ancient building}.
Large-language-model paraphrasing \cite{brown2020gpt3} resolves the contextual
problem but introduces a different one: with no formal constraint on semantic
drift, an unconstrained model may produce a document that preserves surface
fluency while substantially altering meaning, logical structure, or
domain-specific claims.  Neither approach provides a principled account of
\emph{how much} the resulting document has changed or \emph{which aspects} of
the original have been preserved.

Most NLP papers treat the problem as follows: a document is a string; we
transform the string; we measure whether the transformation was successful.
The ontological status of the document is never questioned.  We take a
fundamentally different position.  A document is a point on a semantic
manifold; the text file is merely one coordinate representation of an underlying
geometric object.  Two different strings---\emph{The scientist examined the
structure} and \emph{The researcher analyzed the architecture}---may occupy
nearly the same region of the manifold.  The geometric object is primary;
the text is secondary.

This reframing elevates document perturbation from an NLP engineering task to a
problem in constrained trajectory planning on a Riemannian manifold
\cite{lee2013smooth, docarmo1992riemannian}.  The framework admits a natural
formal language: fiber bundles, metric tensors, connections, holonomy, and
completion-weighted energy allocation.  The central question is no longer
\emph{which words should be replaced} but \emph{which directions in semantic
space are accessible from a given document under a given budget and a given
set of conservation laws}.

We make the following contributions.

\begin{enumerate}[label=(\roman*)]
\item We introduce \emph{semantic equivalence classes} via an approximate
  equivalence relation and show that fibers are exactly these classes,
  grounded in a locally linearized manifold structure (\S\ref{sec:manifold}).
\item We develop the \emph{semantic fiber framework} (\S\ref{sec:fiber}),
  defining semantic sections, reachable sets, and proving that controlled
  augmentation is sampling from a fiber at a specified radius.
\item We introduce \emph{semantic velocity} and \emph{curvature-constrained
  trajectories} (\S\ref{sec:budget}), grounding curvature bounds in classical
  differential geometry.
\item We introduce a \emph{semantic connection} $\connD$ on the bundle
  (\S\ref{sec:budget}) and define admissible perturbations as horizontal lifts,
  resolving the non-holonomic failure modes of embedding-distance-only approaches.
\item We formalize conservation via \emph{Noether-style semantic invariants}
  and \emph{semantic gauge freedom} (\S\ref{sec:conservation}), characterizing
  the admissible perturbation region as a fiber sub-manifold.
\item We prove a \emph{Budget Allocation Principle} for \emph{semantic colored
  noise} (\S\ref{sec:colorednoise}) and introduce \emph{semantic accessibility
  entropy} with a monotonicity theorem.
\item We model mutation profiles as \emph{profile-dependent metric tensors and
  connection structures} (\S\ref{sec:profiles}), deriving a \emph{domain
  rigidity} measure from the ratio of connection curvature to admissible
  trajectory curvature.
\end{enumerate}


% ════════════════════════════════════════════════════════════════════════════
\section{Background}
\label{sec:background}
% ════════════════════════════════════════════════════════════════════════════

\subsection{Document Augmentation}

Synonym substitution using lexical resources such as WordNet
\cite{miller1995wordnet} was among the first systematic approaches to document
augmentation.  Back-translation \cite{sennrich2016backtranslation} introduced
a noise-injection mechanism via round-tripping through an intermediate language.
Large pre-trained language models \cite{raffel2020t5, brown2020gpt3} enable
fluent paraphrase generation but without quantitative control over semantic
drift.  Automatic evaluation metrics such as BERTScore \cite{zhang2020bertscore}
and BLEURT \cite{sellam2020bleurt} provide post-hoc assessment but do not
constrain generation.

\subsection{Semantic Embedding Spaces}

Dense document representations \cite{mikolov2013efficient, devlin2019bert,
reimers2019sentencebert} map text into high-dimensional spaces where cosine
similarity approximates semantic relatedness.  A substantial literature
\cite{ethayarajh2019contextual, rahaman2019spectralbias} has characterized
the geometric structure of these spaces.  Information geometry
\cite{amari2016information} provides a rigorous framework for treating such
spaces as Riemannian manifolds equipped with the Fisher--Rao metric; our
profile-dependent metric tensors are a discrete analogue.

\subsection{Diffusion Processes and Colored Noise}

Denoising diffusion probabilistic models \cite{ho2020ddpm} frame generation as
iterative noise removal along a stochastic trajectory.  The spectral bias of
these processes---coarse structure resolves before fine detail
\cite{rahaman2019spectralbias, falck2025fourierspace}---motivates
\emph{colored noise diffusion sampling} \cite{davidson2026colorednoisediffusionsampling},
which routes perturbation energy toward unresolved frequency bands via a
completion field.  Blue-noise approaches \cite{huang2024bluenoise} and geometric
analyses of sampling trajectories \cite{wang2025rotationaltrajectories} further
develop the idea that perturbation energy should respect existing structure.

\subsection{Constraint-Based Optimization and Control Theory}

Constrained text generation has been approached through lexical constraints
\cite{hokamp2017lexically}, Lagrangian relaxation, and retrieval-augmented
consistency enforcement.  In control theory, constrained trajectory planning
on smooth manifolds \cite{lee2013smooth, docarmo1992riemannian} is a mature
field; our budget constraints and curvature bounds are direct analogues.
Optimal control theory \cite{liberzon2012optimal, bertsekas1995dynamic} provides
the variational framework underlying our budget allocation principle.  Gauge
theories in physics \cite{docarmo1992riemannian} supply the connection and
holonomy language adopted in \S\ref{sec:budget} and \S\ref{sec:conservation}.


% ════════════════════════════════════════════════════════════════════════════
\section{The Semantic Manifold and Equivalence Classes}
\label{sec:manifold}
% ════════════════════════════════════════════════════════════════════════════

Let $\calX$ denote the space of all syntactically valid documents in a given
natural language.  Under a fixed embedding function $\phi : \calX \to \bR^d$
we work in the embedded space $\phi(\calX)$.

\begin{definition}[Semantic embedding]
\label{def:embedding}
A \emph{semantic embedding} is a map $\phi : \calX \to \bR^d$ such that
$\norm{\phi(D)} = 1$ for all $D \in \calX$ and $\phi(D_1) \approx \phi(D_2)$
whenever $D_1$ and $D_2$ are judged semantically equivalent by human annotators.
\end{definition}

\begin{definition}[Semantic manifold]
\label{def:manifold}
The \emph{semantic manifold} $\calM \subset \bR^d$ is the image $\phi(\calX)$,
endowed with the metric inherited from $\bR^d$ restricted to $\calM$.
\end{definition}

\subsection{Local Linearization}

Although $\calM$ may possess complicated global topology, we assume that every
point $m \in \calM$ admits a neighbourhood $U_m$ locally diffeomorphic to
$\bR^k$ for some $k \le d$ \cite{lee2013smooth}.  Consequently, local
perturbation trajectories may be analyzed within the tangent space $T_m\calM$,
a $k$-dimensional real vector space.  This justifies the use of local velocity
vectors, curvature estimates, and metric tensors throughout the paper: all
such objects are defined tangentially and then patched together across
overlapping charts via standard manifold machinery \cite{docarmo1992riemannian}.

\subsection{Tangent Directions as Semantic Operations}

A tangent vector $v \in T_m\calM$ may be interpreted as an infinitesimal
semantic operation: the direction in which the document's meaning moves under
an infinitesimal perturbation.  Concretely, lexical substitutions induce
directions that vary primarily within register and denotational axes; syntactic
transformations perturb argument-structure and information-structure axes;
discourse-level rewrites induce larger displacements spanning multiple semantic
dimensions simultaneously.  These qualitative distinctions motivate a
decomposition of the tangent space into semantically distinguished subspaces:
\[
   T_m\calM
   \;=\;
   T_m^{\mathrm{lex}}
   \oplus
   T_m^{\mathrm{syn}}
   \oplus
   T_m^{\mathrm{sem}},
\]
corresponding to lexical, syntactic, and semantic modes of variation.  The
decomposition need not be orthogonal, and the subspaces may overlap in practice;
nevertheless, the three operator families introduced in \S\ref{sec:budget}
(lexical, syntactic, semantic) are designed to generate motion primarily within
$T_m^{\mathrm{lex}}$, $T_m^{\mathrm{syn}}$, and $T_m^{\mathrm{sem}}$
respectively.

The semantic budget vector $\mathcal{B} = (B_\ell, B_s, B_p, \delta)$
therefore constrains not only the total magnitude of motion in $\calM$ but also
its decomposition across semantically distinct tangent directions.  A budget
that exhausts $B_\ell$ before reaching $B_s$ or $B_p$ has committed the
perturbation to lexical variation and foreclosed the syntactic and semantic modes.
This makes the budget vector a coarse-grained control over the \emph{direction}
of semantic travel, not merely its distance.

\subsection{Semantic Equivalence Classes}

\begin{definition}[Approximate semantic equivalence]
\label{def:equiv}
Two documents $D_1, D_2 \in \calX$ are \emph{semantically $\ve$-equivalent},
written $D_1 \sim_\ve D_2$, if
$\dM\bigl(\phi(D_1),\, \phi(D_2)\bigr) \le \ve$.
\end{definition}

\begin{proposition}[$\sim_\ve$ is an approximate equivalence relation]
\label{prop:equiv}
The relation $\sim_\ve$ is reflexive and symmetric, and transitive up to $2\ve$:
$D_1 \sim_\ve D_2$ and $D_2 \sim_\ve D_3$ imply $D_1 \sim_{2\ve} D_3$.
\end{proposition}

\begin{proof}
Reflexivity and symmetry are immediate.  Transitivity up to $2\ve$ follows from
the triangle inequality: $\dM(\phi(D_1), \phi(D_3)) \le \dM(\phi(D_1),
\phi(D_2)) + \dM(\phi(D_2), \phi(D_3)) \le 2\ve$.
\end{proof}

The equivalence classes under $\sim_\ve$ are the $\ve$-neighbourhoods of each
point on $\calM$; documents in the same class are acceptable paraphrases of one
another at tolerance $\ve$.

\begin{definition}[Semantic neighbourhood]
\label{def:neighbourhood}
$\mathcal{N}_r(D) = \{\, D' \in \calX \mid \dM(\phi(D),\, \phi(D')) \le r\,\}$.
\end{definition}

This establishes the logical chain:
embeddings $\to$ manifold $\to$ local linearization $\to$ tangent decomposition
$\to$ equivalence relation $\to$ neighbourhoods $\to$ fibers $\to$ perturbation.


% ════════════════════════════════════════════════════════════════════════════
\section{Document Perturbation as Fiber Exploration}
\label{sec:fiber}
% ════════════════════════════════════════════════════════════════════════════

\begin{definition}[Semantic projection]
\label{def:projection}
The \emph{semantic projection} is the surjective map
$\vpi : \calX \to \calM$, $\vpi(D) = \phi(D)$.
\end{definition}

\begin{definition}[Semantic fiber]
\label{def:fiber}
The \emph{semantic fiber} over $m \in \calM$ is
$\fiber{m} = \{\, D \in \calX \mid \vpi(D) = m\,\}$.
The \emph{approximate fiber} at tolerance $\ve$ is
$\fiber{m}_\ve = \{D \in \calX : \dM(\vpi(D), m) \le \ve\}$.
\end{definition}

\begin{definition}[Semantic section]
\label{def:section}
A \emph{semantic section} is a map $s : \calM \to \calX$ such that
$\vpi \circ s = \mathrm{id}_{\calM}$.  A section chooses one canonical textual
realization for each semantic identity.
\end{definition}

The original document $D_0$ is therefore not a privileged object but a
\emph{choice of section}: it selects one representative from the fiber
$\fiber{\vpi(D_0)}$.  Controlled augmentation explores other elements of
the same fiber.  Different choices of section over $\calM$---different
canonical representatives---yield different augmentation starting points
for the same underlying semantic content.

\begin{proposition}[Fiber exploration as constrained augmentation]
\label{prop:fiber}
Augmentation within semantic budget $B$ is equivalent to sampling from
$\fiber{\vpi(D_0)}_B$.
\end{proposition}

\begin{proof}
By Definition \ref{def:neighbourhood}, sampling within budget $B$ is sampling
from $\mathcal{N}_B(D_0)$.  By Definition \ref{def:fiber}, the approximate fiber
$\fiber{\vpi(D_0)}_B$ consists of all $D'$ with
$\dM(\phi(D'), \phi(D_0)) \le B$.  Since $\vpi = \phi$, these two sets coincide.
\end{proof}

\subsection{Reachable Sets}

\begin{definition}[Reachable set]
\label{def:reachable}
The \emph{reachable set} of $D$ under budget $B$ is
$\calR_B(D) = \{\, D' \in \calX \mid \dM(\phi(D), \phi(D')) \le B\,\}$.
\end{definition}

A perturbation engine is more \emph{expressive} if it can reach a larger
fraction of $\calR_B(D)$ \cite{liberzon2012optimal}.  By Proposition
\ref{prop:fiber}, $\calR_B(D_0) = \fiber{\vpi(D_0)}_B$: the reachable set
is the approximate fiber.  Budget and fiber radius are the same quantity
measured from different perspectives.

\begin{remark}[Why synonym replacement can fail]
A sequence of individually small steps can exit $\calR_B(D_0)$ if the
trajectory curves sharply.  Two steps of size $\delta$ in perpendicular
directions reach distance $\delta\sqrt{2}$ from the origin; two steps in
opposite directions return to the origin.  Bounding total displacement alone
does not prevent this.
\end{remark}


% ════════════════════════════════════════════════════════════════════════════
\section{Semantic Distance Budgets, Velocity, and the Semantic Connection}
\label{sec:budget}
% ════════════════════════════════════════════════════════════════════════════

\begin{definition}[Perturbation operator]
\label{def:operator}
A \emph{perturbation operator} $\calO : \calX \to \calX$ maps a document to a
modified version.  Operators are organized by scale:
\begin{enumerate}[label=(\roman*)]
  \item \emph{Lexical operators}: synonym substitution, register shifts.
        Typical cost $\Delta \approx 0.001$--$0.012$.
  \item \emph{Syntactic operators}: active--passive, clause reordering.
        Typical cost $\Delta \approx 0.008$--$0.020$.
  \item \emph{Semantic operators}: LLM paraphrase, summarization--expansion.
        Typical cost $\Delta \approx 0.015$--$0.040$.
\end{enumerate}
The \emph{cost} of applying $\calO$ to $D$ is
$c(\calO, D) = \dM(\phi(D), \phi(\calO(D)))$.
\end{definition}

\begin{definition}[Semantic budget vector]
\label{def:budget}
A \emph{semantic budget vector} is
$\mathcal{B} = (B_\ell, B_s, B_p, \delta) \in \bR_{>0}^4$
specifying lexical, syntactic, and semantic budgets and a maximum step size
$\delta > 0$.  The total budget is $B = B_\ell + B_s + B_p$.
\end{definition}

\begin{definition}[Admissible perturbation sequence]
\label{def:admissible}
A sequence of operators producing trajectory $D_0 \to D_1 \to \cdots \to D_n$
is \emph{admissible} with respect to $\mathcal{B}$ if:
\[
   \dM\bigl(\phi(D_0),\, \phi(D_n)\bigr) \;\le\; B
   \quad\text{(total displacement),}
\]
\[
   \dM\bigl(\phi(D_{k-1}),\, \phi(D_k)\bigr) \;\le\; \delta
   \quad\forall\, k\quad\text{(step bound).}
\]
\end{definition}

The perturbation problem is:
\begin{equation}
\label{eq:opt}
   \max_{(\calO_1,\ldots,\calO_n)\;\text{admissible}} V(D_n)
\end{equation}
where $V : \calX \to \bR$ is a diversity functional, typically
$V(D') = \dM(\phi(D_0), \phi(D'))$.

\subsection{Semantic Velocity and Curvature}

\begin{definition}[Semantic velocity]
\label{def:velocity}
The \emph{semantic velocity} at step $k$ is
$v_k = \phi(D_k) - \phi(D_{k-1}) \in T_{\phi(D_{k-1})}\calM \subset \bR^d$.
\end{definition}

\begin{definition}[Discrete trajectory curvature]
\label{def:curvature}
The \emph{discrete curvature} at step $k$ is
\[
   \kappa_k \;=\; 1 - \frac{\inner{v_k}{v_{k+1}}}{\norm{v_k}\,\norm{v_{k+1}}}.
\]
$\kappa_k = 0$ is straight-line motion; $\kappa_k = 2$ is reversal.
\end{definition}

A \emph{curvature-constrained trajectory} satisfies $\kappa_k \le \kappa_{\max}$
at all interior steps.

\begin{proposition}[Curvature reduces effective reachable radius]
\label{prop:curvature}
For $n = B/\delta$ maximum-size steps, bounding $\kappa_k \le \kappa_{\max} < 1$
reduces maximum possible displacement below the unconstrained bound $\sqrt{n}\,\delta$.
\end{proposition}

\begin{proof}[Proof sketch]
Unconstrained orthogonal steps yield $\norm{\sum_k v_k} \le \sqrt{n}\,\delta$
by Cauchy--Schwarz.  With $\kappa_k \le \kappa_{\max}$, consecutive velocities
lie within a cone of half-angle $\theta = \arccos(1-\kappa_{\max})$.  By
induction, the $n$-step displacement is bounded by
$\delta \sum_{k=0}^{n-1} \cos^k\!\theta < \sqrt{n}\,\delta$ for $\kappa_{\max} < 1$.
\end{proof}

\subsection{Parallel Transport of Meaning and the Semantic Connection}

The distance metric $g$ specifies the cost of movement but does not specify
how semantic quantities are transported along a trajectory.  Certain
semantically catastrophic mutations---negation reversal, causal inversion---are
locally distance-preserving in embedding space while producing globally
inadmissible documents.  To address this failure mode we introduce a connection
on the semantic bundle.

\begin{definition}[Semantic connection]
\label{def:connection}
A \emph{semantic connection} $\connD$ on the bundle $\vpi : \calX \to \calM$
is a rule specifying, for each trajectory $D_0 \to D_1 \to \cdots \to D_n$,
how semantic quantities are transported from fiber to fiber.  A quantity $Q$
is \emph{parallel transported} along the trajectory when
\[
   \connD_{\dot D} Q \;=\; 0.
\]
\end{definition}

Admissible perturbations correspond to \emph{horizontal trajectories} of the
bundle: trajectories whose velocity lies entirely in the horizontal distribution
$H \subset T\calX$ defined by $\connD$.  Inadmissible semantic alterations
correspond to motion with a nonzero vertical component, moving between fibers
rather than within them.

\subsection{Holonomy and Accumulated Drift}

A perturbation trajectory may return to a neighbourhood of its starting point
while failing to preserve all transported semantic quantities.  Let $\Gamma$
be a closed perturbation loop in $\calM$---a sequence of operators that returns
the document's embedding to its original position.  Parallel transport around
$\Gamma$ via $\connD$ induces a \emph{holonomy operator}
\[
   \mathrm{Hol}_{\Gamma} : \fiber{\vpi(D_0)} \to \fiber{\vpi(D_0)}.
\]
If $\mathrm{Hol}_{\Gamma} \neq \mathrm{id}$, then even a closed embedding
trajectory accumulates semantic drift: the document returns to the same
embedding position but ends in a different semantic orbit, with some conserved
quantity $Q$ shifted by the holonomy.

Holonomy is nonzero precisely when the curvature $\curv = d\connD + \connD
\wedge \connD$ of the connection is nontrivial over the region enclosed by
$\Gamma$.  This provides a direct operational interpretation of the curvature
tensor: $\curv_P$ measures how readily long perturbation chains accumulate
invisible semantic drift even when each individual step is admissible.  A
high-curvature connection means that closed loops are dangerous; a flat
connection means that order of operations is irrelevant and semantic drift
cannot accumulate.  The domain rigidity measure $R(P)$ defined in
\S\ref{sec:profiles} formalizes the ratio of how much holonomy the connection
can accumulate to how much the discourse community will tolerate.

\begin{example}[Negation as a vertical direction]
Introduce a semantic polarity coordinate $P(D) \in \{+1, -1\}$ where $+1$
denotes affirmation and $-1$ denotes negation.  Define $\connD$ so that
$P$ is a horizontally conserved quantity: $\connD_{\dot D} P = 0$ along any
admissible trajectory.  A mutation changing $P$ is not merely expensive; it
moves in a vertical direction and exits the horizontal distribution entirely.
The trajectory is inadmissible regardless of its embedding-space cost.
\end{example}

\begin{example}[Causality as graph conservation]
Let $G_D = (V, E)$ be the directed causal graph of document $D$, where nodes
are propositions and edges encode causal or inferential dependencies.  Define
$\connD$ so that $G_D$ is horizontally conserved: admissible perturbations
must preserve $G_D$ up to graph isomorphism.  The documents
\emph{A caused B} and \emph{B caused A} are not connected by any horizontal
path in the bundle, even though their embedding vectors may be close, because
any path between them must pass through a vertical displacement that inverts
the edge direction.
\end{example}

The connection generalizes the projection-based conservation constraints of
\S\ref{sec:conservation} by making conservation \emph{path-dependent} rather
than purely position-dependent.  A quantity conserved via the connection remains
invariant under parallel transport along any admissible trajectory; a
position-based constraint merely requires that the quantity be the same at the
endpoint.  This is precisely the distinction between holonomic and non-holonomic
constraints in classical mechanics: the current conservation formalism is
holonomic (positions constrained), while the connection is non-holonomic
(velocities/directions constrained).


% ════════════════════════════════════════════════════════════════════════════
\section{Conservation Laws}
\label{sec:conservation}
% ════════════════════════════════════════════════════════════════════════════

\begin{definition}[Protected concept set]
\label{def:protected}
A \emph{protected concept set} is
$\mathcal{C} = \{c_1, \ldots, c_k\}$, each $c_i$ represented by
$\phi(c_i) \in \bR^d$.
\end{definition}

\begin{definition}[Conservation constraint]
\label{def:conservation}
A perturbation $D \to D'$ satisfies the \emph{conservation constraint} with
respect to $\mathcal{C}$ and tolerance $\ve > 0$ if
\[
   \bigl|\inner{\phi(D')}{\phi(c_i)} - \inner{\phi(D)}{\phi(c_i)}\bigr|
   \;\le\; \ve
   \qquad \forall\, c_i \in \mathcal{C}.
\]
\end{definition}

\subsection{Noether-Style Semantic Invariants}

\begin{definition}[Semantic conserved quantity]
\label{def:noether}
A function $Q : \calX \to \bR$ is a \emph{semantic conserved quantity} if
every admissible perturbation $D \to D'$ satisfies
$|Q(D') - Q(D)| \le \ve$ for some declared tolerance $\ve$.
\end{definition}

\begin{example}[Conserved charges in practice]
\begin{enumerate}[label=(\roman*)]
  \item \emph{Thesis projection}: $Q(D) = \inner{\phi(D)}{\phi(\text{central claim})}$.
  \item \emph{Terminology}: $Q(D) = \inner{\phi(D)}{\phi(t)}$ for each
        domain term $t \in \{\text{RSVP}, \text{TARTAN}, \text{Friedmann}, \ldots\}$.
  \item \emph{Citation preservation}: fraction of source citations appearing in $D$.
  \item \emph{Numerical claim conservation}: set of quantitative assertions
        in $D$, measured by entailment score.
  \item \emph{Polarity}: $Q(D) = P(D) \in \{+1, -1\}$, conserved via $\connD$.
  \item \emph{Causal graph}: $Q(D) = G_D$ up to isomorphism, conserved via $\connD$.
\end{enumerate}
\end{example}

Each conserved quantity defines a constraint hyper-surface in $\calM$.

\begin{proposition}[Invariant fiber sub-manifold]
\label{prop:invariant}
Let $Q_1, \ldots, Q_n$ be smooth semantic conserved quantities with values
$q_i = Q_i(D_0)$.  Then the admissible perturbation space
\[
   \calF_Q \;=\; \bigl\{\, D \in \calX \;\mid\; Q_i(D) = q_i,\;
   i = 1, \ldots, n \,\bigr\}
\]
forms a sub-manifold of the semantic fiber $\fiber{\vpi(D_0)}$, provided the
gradients $\nabla Q_i$ are linearly independent on $\calF_Q$.
\end{proposition}

\begin{proof}
By the regular level set theorem \cite{lee2013smooth}, the preimage of a
regular value of a smooth map $Q = (Q_1, \ldots, Q_n) : \calX \to \bR^n$
is a smooth sub-manifold of $\calX$ of codimension $n$.  Since $\calF_Q
\subset \fiber{\vpi(D_0)}$ by construction, it is a sub-manifold of the fiber.
\end{proof}

Conservation is therefore not merely a constraint; it is geometry.  The
perturbation engine explores a lower-dimensional sub-manifold of the fiber,
one whose dimension is reduced by one for each independent conserved quantity.

\subsection{Semantic Gauge Freedom}

Not every modification represents genuine semantic change.  Certain surface
variations alter textual realization while preserving all conserved quantities.

\begin{definition}[Semantic gauge transformation]
\label{def:gauge}
A \emph{semantic gauge transformation} is a map $g : \calX \to \calX$ such that
$Q_i(g(D)) = Q_i(D)$ for all conserved quantities $Q_i$ and all $D \in \calX$.
Documents related by a gauge transformation are \emph{gauge equivalent}.
\end{definition}

\begin{example}[Gauge-equivalent substitutions]
The substitutions \emph{researcher} $\leftrightarrow$ \emph{scientist} and
\emph{investigate} $\leftrightarrow$ \emph{analyze} are gauge transformations:
they alter surface realization without moving the document's projection onto
thesis, polarity, or causal graph axes.
\end{example}

The perturbation problem therefore has two levels: finding gauge-equivalent
realizations (cheap, within the same orbit of the gauge group), and genuinely
exploring the fiber by moving between gauge orbits (more expensive, consumes
semantic budget).  The semantic budget $B$ governs inter-orbit travel; gauge
transformations are free with respect to the budget.


% ════════════════════════════════════════════════════════════════════════════
\section{Semantic Colored Noise and Accessibility Entropy}
\label{sec:colorednoise}
% ════════════════════════════════════════════════════════════════════════════

Standard document perturbation is implicitly \emph{semantic white noise}: it
allocates perturbation energy uniformly across all semantic components,
regardless of how much each has already been varied.  We show that this is
suboptimal and introduce a completion-weighted allocation principle motivated
by colored noise diffusion sampling
\cite{davidson2026colorednoisediffusionsampling}.

\subsection{Semantic Completion Fields}

Let $\mathcal{C}_{\mathrm{sem}} = \{c_1, \ldots, c_m\}$ be a decomposition of
the document's semantic content into components, concretely:
$\{\text{thesis},\; \text{evidence},\; \text{terminology},\;
  \text{style},\; \text{narrative},\; \text{examples}\}$.

\begin{definition}[Semantic completion field]
\label{def:completion}
The \emph{semantic completion field} is a function
$\gamma : \mathcal{C}_{\mathrm{sem}} \times [0,1] \to [0,1]$
where $\gamma(c, t)$ measures the fraction of variation budget already spent on
component $c$ by perturbation step $t$.
\end{definition}

\begin{example}[Completion landscape of a mature document]
A document with a well-developed thesis and established terminology but
relatively unvaried style and few examples might exhibit:
$\gamma(\text{thesis}) = 0.95$, $\gamma(\text{terminology}) = 0.92$,
$\gamma(\text{style}) = 0.40$, $\gamma(\text{examples}) = 0.25$.
A uniform perturbation engine wastes budget on thesis and terminology;
a completion-weighted engine routes variation toward style and examples.
\end{example}

\subsection{Semantic Utility and the Budget Allocation Principle}

\begin{definition}[Semantic utility function]
\label{def:utility}
The \emph{semantic utility} of allocating perturbation energy $E$ to component
$c$ with current completion $\gamma(c)$ is
\[
   U_c(E) \;=\; \sqrt{E\,(1 - \gamma(c))},
\]
representing diminishing returns in highly completed semantic regions: additional
budget applied to a saturated component yields progressively less new variation.
\end{definition}

\begin{definition}[Semantic colored noise]
\label{def:colorednoise}
A perturbation strategy produces \emph{semantic colored noise} if energy
allocated to component $c$ at step $t$ is
\[
   E(c, t) \;=\; \mathcal{B} \cdot
   \frac{1 - \gamma(c, t)}{\displaystyle\sum_j \bigl(1 - \gamma(j, t)\bigr)}.
\]
Uniform allocation $E(c, t) = \mathcal{B}/m$ is \emph{semantic white noise}.
\end{definition}

\begin{theorem}[Budget Allocation Principle]
\label{thm:allocation}
Let $V(D') = \sum_c \sigma_c \cdot U_c(E(c))$ be the total semantic utility
with weights $\sigma_c > 0$ and utility $U_c$ as in Definition \ref{def:utility}.
Under the constraint $\sum_c E(c) = \mathcal{B}$ with $E(c) \ge 0$, semantic
colored noise allocation maximizes $V$ over all budget distributions.
\end{theorem}

\begin{proof}
Maximize $\sum_c \sigma_c \sqrt{E(c)(1-\gamma(c))}$ subject to
$\sum_c E(c) = \mathcal{B}$.  By the method of Lagrange multipliers, at the
optimum $\frac{\partial}{\partial E(c)} \bigl[\sigma_c \sqrt{E(c)(1-\gamma(c))}
- \lambda E(c)\bigr] = 0$, giving
$\sigma_c \sqrt{(1-\gamma(c))} / (2\sqrt{E^*(c)}) = \lambda$, hence
$E^*(c) \propto \sigma_c^2(1-\gamma(c))$.  With uniform weights $\sigma_c = 1$
this yields $E^*(c) \propto 1-\gamma(c)$, which is exactly the colored noise
allocation of Definition \ref{def:colorednoise}.
\end{proof}

\begin{remark}[Relation to CNS]
Davidson et al.\ \cite{davidson2026colorednoisediffusionsampling} route
diffusion energy toward $\sqrt{1-\gamma(f,t)}$ over Fourier frequency bands.
Theorem \ref{thm:allocation} establishes that the same energy-routing principle
is optimal in the semantic domain when frequency bands are replaced by semantic
components.  The mathematical structure is identical; what changes is the space
over which completion is measured.
\end{remark}

\begin{remark}[Semantic spectral bias]
The CNS paper \cite{davidson2026colorednoisediffusionsampling} identifies a
\emph{spectral bias} in diffusion models: coarse structure resolves before fine
structure.  The semantic analogue is \emph{conceptual bias}: core thesis and
terminology are established early and become increasingly rigid, while style,
examples, and phrasing remain underexplored.  Theorem \ref{thm:allocation}
shows that completion-weighted allocation is not merely heuristically appealing
but provably optimal under the utility model of Definition \ref{def:utility}.
\end{remark}

\subsection{Semantic Accessibility Entropy}

\begin{definition}[Semantic accessibility entropy]
\label{def:entropy}
The \emph{semantic accessibility entropy} of a document $D$ under budget $B$
and active profile metric $g_P$ is
\[
   S(D, B) \;=\; \log \operatorname{Vol}_P\!\bigl(\calR_B(D)\bigr),
\]
where $\operatorname{Vol}_P$ denotes volume measured under the profile metric.
\end{definition}

This quantity measures the logarithmic volume of semantic space accessible from
$D$ while respecting all budget and conservation constraints.  Increasing budget
increases semantic accessibility; conservation laws reduce it by restricting the
admissible region to the sub-manifold $\calF_Q \subset \fiber{\vpi(D)}$.

\begin{proposition}[Monotonicity of semantic accessibility]
\label{prop:entropy}
For $B_1 < B_2$, we have $S(D, B_1) \le S(D, B_2)$.
\end{proposition}

\begin{proof}
Since $B_1 < B_2$ implies $\calR_{B_1}(D) \subseteq \calR_{B_2}(D)$,
their volumes satisfy
$\operatorname{Vol}_P(\calR_{B_1}(D)) \le \operatorname{Vol}_P(\calR_{B_2}(D))$,
and taking logarithms preserves the inequality.
\end{proof}

\subsection{Entropy Reduction by Conservation}

Conservation laws do not merely filter out inadmissible perturbations after the
fact; they geometrically reduce the volume of the accessible region and therefore
strictly decrease semantic accessibility entropy.

\begin{proposition}[Conservation reduces accessibility]
\label{prop:conservation_entropy}
Let $\calF_Q \subseteq \calR_B(D)$ be the admissible region determined by a
collection of conserved quantities $Q_1, \ldots, Q_n$ as in Proposition
\ref{prop:invariant}.  Define the \emph{constrained accessibility entropy}
$S_Q(D, B) = \log \operatorname{Vol}_P(\calF_Q)$.  Then
\[
   \operatorname{Vol}_P(\calF_Q) \;\le\; \operatorname{Vol}_P\!\bigl(\calR_B(D)\bigr),
\]
and consequently $S_Q(D, B) \le S(D, B)$.
\end{proposition}

\begin{proof}
Since $\calF_Q \subseteq \calR_B(D)$ by construction, monotonicity of
$\operatorname{Vol}_P$ under set inclusion gives
$\operatorname{Vol}_P(\calF_Q) \le \operatorname{Vol}_P(\calR_B(D))$.
Taking logarithms preserves the inequality, yielding $S_Q(D,B) \le S(D,B)$.
\end{proof}

Each independent conserved quantity reduces the dimension of $\calF_Q$ by one
(Proposition \ref{prop:invariant}), so the entropy reduction is not merely
bounded but progressive: adding conserved quantities strictly decreases $S_Q$.
The semantic connection $\connD$ enforces further reduction beyond what
position-based conserved quantities alone can achieve, because horizontal
trajectories must avoid not just forbidden endpoints but forbidden
\emph{directions} at every step.  In RSVP terms \cite{flyxion2025rsvp},
$S(D, B)$ is the semantic entropy of the document: not a measure of disorder
but of the volume of admissible futures.  The Budget Allocation Principle is
the optimal strategy for allocating that volume across semantic dimensions.


% ════════════════════════════════════════════════════════════════════════════
\section{Profile-Dependent Metric Tensors and Domain Rigidity}
\label{sec:profiles}
% ════════════════════════════════════════════════════════════════════════════

Different discourse communities define semantic equivalence differently.  An
academic paraphrase may freely substitute Latinate for Germanic vocabulary
without semantic loss; a legal paraphrase of the same sentence may not, because
the choice between \emph{shall} and \emph{must} carries distinct deontic force.
These differences reflect genuinely different metric structures on $\calM$ and
different connection structures on the bundle.

\begin{definition}[Profile geometry]
\label{def:profile_geometry}
A \emph{profile geometry} for community $P$ is a triple $(g_P, \connD_P, \curv_P)$
where $g_P$ is a Riemannian metric on $\calM$ (encoded as a matrix
$M_P \in \bR^{d\times d}$, $M_P \succ 0$), $\connD_P$ is a semantic connection
compatible with $g_P$, and $\curv_P = d\connD_P + \connD_P \wedge \connD_P$
is the curvature of that connection.  The \emph{profile distance} is
$d_P(a,b) = \sqrt{(a-b)^\top M_P (a-b)}$.
\end{definition}

\begin{example}[Academic metric]
Let $\hat{u}_{\mathrm{form}}$ be the unit vector from casual to formal register
in embedding space.  Then
$M_{\mathrm{acad}} = I + (w_{\mathrm{form}}-1)\hat{u}_{\mathrm{form}}
\hat{u}_{\mathrm{form}}^\top$, $w_{\mathrm{form}} \gg 1$,
makes register-shifting moves expensive under $d_{\mathrm{acad}}$.
\end{example}

\begin{example}[Legal metric]
Under a legal profile, the deontic modality axis (obligation, permission,
prohibition) carries amplified weight in $M_{\mathrm{legal}}$.  The connection
$\connD_{\mathrm{legal}}$ forbids horizontal trajectories that cross the
obligation--permission boundary, so substitutions altering deontic force are
inadmissible regardless of their embedding distance.
\end{example}

\begin{proposition}[Profile independence of the perturbation algorithm]
\label{prop:profile_independence}
The perturbation engine solving \eqref{eq:opt} is profile-independent: it calls
$d_P(\cdot,\cdot)$ and $\connD_P$ as oracles and runs identically for all
profiles.  Changing profile changes the geometry of the search without altering
the algorithm.
\end{proposition}

\subsection{Metric Learning Interpretation}

The profile metric $M_P$ is a \emph{learned Mahalanobis metric}
\cite{amari2016information}.  Given pairs $(D_i, D_j)$ labeled equivalent
under profile $P$ by domain experts, $M_P$ is learned by minimizing
\[
   \sum_{\text{equiv}} d_P(\phi(D_i),\phi(D_j))^2 +
   \lambda \sum_{\text{distinct}} \max\!\bigl(0,\, \alpha - d_P(\phi(D_i),\phi(D_j))\bigr)
\]
subject to $M_P \succeq 0$.  This connects profile metrics directly to the
metric-learning literature \cite{larsen2016autoencoding}.  The connection
$\connD_P$ requires richer supervision: pairs of documents together with
\emph{trajectories} between them, labeled as admissible or inadmissible by
domain experts, are needed to learn the horizontal distribution.

\subsection{Curvature and Domain Rigidity}

\begin{definition}[Domain rigidity]
\label{def:rigidity}
Let $\curv_P$ denote the curvature of the semantic connection under profile $P$
and $\kappa_{\max,P}$ the maximum admissible trajectory curvature.  The
\emph{rigidity} of discourse community $P$ is
\[
   R(P) \;=\; \frac{\norm{\curv_P}}{\kappa_{\max,P}}.
\]
\end{definition}

High rigidity characterizes domains where the semantic connection is highly
curved (parallel transport is strongly path-dependent) and the discourse
community tolerates little trajectory curvature.  Low rigidity characterizes
domains where the connection is flatter and larger loops are permitted without
identity failure.

\begin{example}[Rigidity across profiles]
Legal, mathematical, and formal-logical language are expected to exhibit high
rigidity: $\norm{\curv_{\mathrm{legal}}} \gg \norm{\curv_{\mathrm{literary}}}$
while $\kappa_{\max,\mathrm{legal}} \ll \kappa_{\max,\mathrm{literary}}$.
In a legal document, the loop
\emph{obligation} $\to$ \emph{permission} $\to$ \emph{exception} $\to$
\emph{obligation}
may accumulate holonomy: the document returns lexically close to its origin
but legally altered.  In literary prose, an analogous loop through tonal
registers may be perfectly admissible.
\end{example}

The rigidity measure gives a practical implementation rule: a perturbation
engine should use profile-dependent gauges $(g_P, \connD_P, \curv_P)$ rather
than a single universal semantic budget.  For legal text, most admissible motion
is horizontal and narrow.  For literary text, the horizontal distribution is
wider and the system can tolerate larger loops without identity failure.

\begin{remark}[Holonomy as semantic drift]
In gauge theory, non-trivial holonomy (a loop in the base space that fails to
close in the total space) signals the presence of curvature.  In the semantic
bundle, holonomy corresponds to meaning drift accumulated along a closed
trajectory: a sequence of perturbations that returns to the same embedding
position but leaves the document in a different semantic orbit.  High-curvature
connections ($\norm{\curv_P}$ large) make holonomy more likely and more
consequential.  Rigidity $R(P)$ measures the ratio of how much holonomy the
connection can accumulate to how much the discourse community will tolerate.
\end{remark}


% ════════════════════════════════════════════════════════════════════════════
\section{Multi-Scale Perturbation Operators}
\label{sec:operators}
% ════════════════════════════════════════════════════════════════════════════

Each operator scale draws from its own budget axis and must satisfy both the
metric-based budget constraint and the connection-based horizontality condition.

\subsection{Lexical Operators}

Synonym substitution uses WordNet \cite{miller1995wordnet} with
frequency-weighted selection.  The weight exponent
$w_i = \mathrm{freq}(s_i)^{\alpha}$, $\alpha = 1 - 0.85\,t$,
interpolates between frequency-dominated selection at $t=0$ and near-uniform
sampling at $t=1$.  The per-token replacement probability $p(t) = 0.01 + 0.08\,t$
controls substitution frequency independently of synonym pool width.  At
$t \ge 0.5$, an LLM backend is invoked first; WordNet serves as fallback.
All proposed substitutions are checked against the conservation constraint and
the connection's horizontal condition before commitment.

\subsection{Syntactic Operators}

Syntactic operators restructure sentences: active--passive conversion, clause
reordering, appositive insertion.  Syntactic operators incur higher semantic
cost ($\Delta \approx 0.008$--$0.020$) and are particularly susceptible to
connection violations: the documents \emph{A caused B} and \emph{B caused A}
differ syntactically in the same way as \emph{A followed B} and \emph{B followed A},
but only the latter is a safe paraphrase under causal graph conservation.
All syntactic proposals are submitted to the fiber explorer for both metric
and connection checking before commitment.

\subsection{Semantic Operators}

LLM paraphrase operators rewrite at the sentence or paragraph level
($\Delta \approx 0.015$--$0.040$ per sentence).  They are invoked only when
remaining semantic budget $B_p - \mathrm{spent}_p$ exceeds a threshold.  The
LLM system prompt encodes profile information, directing the model toward
register-appropriate outputs; the prompt also includes protected terms and
polarity constraints, partially encoding the connection structure in natural
language.


% ════════════════════════════════════════════════════════════════════════════
\section{Implementation Architecture}
\label{sec:implementation}
% ════════════════════════════════════════════════════════════════════════════

The \textsc{DocPerturb} prototype implements the framework in approximately
3,000 lines of Python and TypeScript across five architectural layers.

The \emph{embedding layer} wraps a sentence-transformer model
\cite{reimers2019sentencebert} and provides sentence-level cached embeddings.
For a document with 200 candidate substitutions, budget checking without
caching requires 200 forward passes; sentence-level caching reduces this to
$n_{\mathrm{sentences}}$ passes by recomputing only the modified sentence.

The \emph{semantic budget tracker} maintains the trajectory $D_0 \to D_k$,
computing total displacement and per-step distance for every proposed mutation.
The completion field $\gamma(c, t)$ is tracked incrementally, updating after
each committed mutation, and budget allocation follows Theorem \ref{thm:allocation}.

The \emph{constraint and connection engine} implements the conservation
constraint (Definition \ref{def:conservation}), the gauge-equivalence check
(Definition \ref{def:gauge}), and a partial implementation of the connection
horizontality condition (negation polarity, simple causal patterns via
dependency parsing).  Full graph-level conservation is reserved for future work.

The \emph{perturbation operators} implement lexical, syntactic, and semantic
mutations, each submitting proposals to a unified \texttt{FiberExplorer.admit()}
interface that checks metric budget, connection horizontality, and conservation
before commitment.

The \emph{search loop} performs greedy optimization of \eqref{eq:opt},
allocating budget to components according to the colored noise principle.
Future work will explore beam search and simulated annealing within the budget
constraint \cite{bertsekas1995dynamic}.


% ════════════════════════════════════════════════════════════════════════════
\section{Experimental Evaluation}
\label{sec:experiments}
% ════════════════════════════════════════════════════════════════════════════

We outline experimental protocols; full empirical results are reserved for a
follow-up study with access to human annotators and standardized benchmark corpora.

\paragraph{Semantic fidelity versus temperature.}
For each temperature $t \in \{0.1, 0.3, 0.5, 0.7, 0.9\}$ and each of the six
profiles, we generate 100 variants of 50 source documents across diverse genres.
We measure embedding cosine distance, human judgments on a 5-point scale, and
preservation of named entities and key claims.

\paragraph{Colored noise versus white noise.}
We compare uniform budget allocation against the completion-weighted allocation
of Theorem \ref{thm:allocation}, measuring semantic diversity (reachable set
hull volume), conservation constraint satisfaction, and human naturalness ratings.

\paragraph{Connection violation detection.}
We measure the rate at which the connection-based horizontality check (negation
and causal graph conservation) blocks mutations that pass the metric-based budget
check but would alter propositional meaning, evaluating the incremental benefit
of connection checking over metric checking alone.

\paragraph{Domain rigidity measurement.}
For each profile, we empirically estimate $R(P) = \norm{\curv_P}/\kappa_{\max,P}$
by measuring how frequently valid (budget-admissible) perturbations produce
meaning changes judged unacceptable by domain experts.  We predict
$R(\mathrm{legal}) \gg R(\mathrm{creative})$.

\paragraph{Retrieval robustness and deduplication benchmarking.}
Perturbed variants at calibrated budgets test whether dense retrieval systems
correctly identify near-duplicates.  Budget $B$ provides the ground-truth
semantic distance label: pairs within $B = 0.05$ are ``equivalent,'' between
$0.05$ and $0.15$ ``related,'' beyond $0.15$ ``distinct.''


% ════════════════════════════════════════════════════════════════════════════
\section{Applications}
\label{sec:applications}
% ════════════════════════════════════════════════════════════════════════════

\emph{Data augmentation.}  Perturbations within small budget $B$ preserve class
labels while increasing diversity; budget provides a formal guarantee that the
label has not been corrupted.

\emph{RAG robustness evaluation.}  Semantic perturbations of queries and source
documents at calibrated distances test retrieval system sensitivity to paraphrase.

\emph{Plagiarism detection calibration.}  Perturbations at increasing budgets
generate test cases at known semantic distances, identifying the budget threshold
at which detection fails.

\emph{Adversarial robustness.}  Selectively violating conservation constraints
or connection horizontality generates adversarial examples targeting specific
semantic structures.

\emph{Theory version control.}  Any object admitting a semantic embedding---a
scientific theory, a software specification, a legal statute---can be tracked
under the framework.  A theory evolves admissibly if each revision stays within
a semantic budget relative to the founding formulation, and conserved quantities
ensure core claims are not inadvertently abandoned.


% ════════════════════════════════════════════════════════════════════════════
\section{Limitations and Open Problems}
\label{sec:limitations}
% ════════════════════════════════════════════════════════════════════════════

\paragraph{Metric tensor and connection construction.}
Profile metrics are currently initialized from hand-specified contrastive pairs;
connections are partially implemented.  Without learned connections, the
horizontal distribution is only approximately enforced via LLM prompting.
Learning $\connD_P$ from trajectory-level supervision is the most important
open problem.

\paragraph{Computational cost.}
Each proposed mutation requires an embedding forward pass and, if connection
checking is active, a dependency parse and graph comparison.  Sentence-level
caching and batched mutation proposals mitigate cost; approximate fiber-membership
testing via locality-sensitive hashing is a further direction.

\paragraph{Completion field estimation.}
Theorem \ref{thm:allocation} assumes $\gamma(c,t)$ can be measured.  In
practice, decomposing a document's semantic content into well-defined components
and estimating completion from embeddings alone is non-trivial.  Structured
semantic representations (knowledge graphs, frame-semantic parses) would make
this reliable.

\paragraph{Non-holonomic language phenomena.}
The connection formalizes the intuition that certain semantic failures are
directional rather than positional.  But the full non-holonomic structure of
language---scope, presupposition, implicature, discourse coherence---is far
richer than the polarity and causality examples developed here.  Mapping the
full horizontal distribution of natural language remains an open problem in
formal semantics.

\paragraph{Embedding faithfulness.}
The framework rests on the assumption that embedding cosine distance approximates
human semantic distance, which is known to fail for negation, presupposition,
and other compositional structures.  The connection addresses some failures;
others require structured representations beyond sentence embeddings.


% ════════════════════════════════════════════════════════════════════════════
\section{Future Directions}
\label{sec:future}
% ════════════════════════════════════════════════════════════════════════════

\emph{Learning the semantic connection.}  Tracking how causal graphs, rhetorical
trees, and polarity structures are preserved across human-authored paraphrases
in large corpora would effectively map the horizontal distribution of the semantic
bundle, making the connection data-driven rather than rule-based.

\emph{Adaptive budget allocation.}  Updating $\gamma(c,t)$ after each committed
mutation and reallocating budget dynamically would better track the actual
exploration frontier than the current static initialization.

\emph{Graph-based conservation.}  Replacing flat embeddings with knowledge-graph
or frame-semantic representations would enable conservation at the propositional
level, addressing negation and causality failure modes within the conservation
formalism rather than via ad-hoc rules.

\emph{Multimodal perturbation.}  The budget, connection, and conservation
framework generalizes to any modality admitting a semantic embedding.
Multi-modal documents could be perturbed jointly under a shared semantic budget
and a joint connection.

\emph{Connection to RSVP and projection formalisms.}  The fiber exploration
framing connects naturally to the RSVP field-theoretic framework
\cite{flyxion2025rsvp}, in which admissible configurations satisfy a projection
condition analogous to fiber membership and entropy measures the volume of
admissible future trajectories.  The semantic accessibility entropy $S(D,B)$
is, in RSVP terms, the semantic entropy of the document.  The Budget Allocation
Principle is the optimal strategy for allocating that entropy across semantic
dimensions.  The connection $\connD$ and its curvature $\curv_P$ correspond to
the field structure of the RSVP plenum, and admissible perturbations correspond
to horizontal lifts of semantic trajectories in the same way that admissible
physical configurations correspond to field configurations satisfying the RSVP
projection conditions.  Diffusion sampling (CNS), document mutation (DocPerturb),
and RSVP accessibility dynamics are therefore all instances of a single abstract
principle: constrained finite-budget perturbation in which energy is routed
toward unresolved dimensions rather than distributed uniformly.


% ════════════════════════════════════════════════════════════════════════════
\section{Conclusion}
\label{sec:conclusion}
% ════════════════════════════════════════════════════════════════════════════

We have argued that document perturbation is fundamentally a geometric problem.
A document is a point on a semantic manifold; the text file is one coordinate
representation of an underlying geometric object; admissible paraphrases occupy
the same semantic fiber; and controlled augmentation is constrained trajectory
planning within that fiber under a budget of semantic distance and a connection
that governs admissible transport of meaning.

The framework contributes a complete mathematical narrative: embedding $\to$
manifold with local linearization $\to$ tangent decomposition $\to$ equivalence
classes $\to$ fibers and sections $\to$ reachable sets $\to$ budget vectors $\to$
semantic velocity and curvature $\to$ semantic connection and horizontal lifts
$\to$ holonomy and accumulated drift $\to$ conservation as geometry
(Noether-style invariants, gauge freedom, invariant sub-manifold) $\to$ entropy
reduction by conservation $\to$ colored noise allocation (Budget Allocation
Principle, accessibility entropy, monotonicity) $\to$ profile-dependent geometry
and domain rigidity.

Each step does genuine mathematical work.  The tangent decomposition connects
the manifold geometry directly to the operator hierarchy.  The semantic
connection resolves the non-holonomic failure modes that distance metrics alone
cannot detect.  The holonomy subsection gives curvature a concrete operational
meaning before it reappears in the profile section.  Proposition
\ref{prop:conservation_entropy} makes precise the verbal claim that conservation
laws reduce semantic accessibility, providing the formal link between the
conservation and entropy threads.  The Budget Allocation Principle shows that
uniform perturbation is suboptimal.  The domain rigidity measure
$R(P) = \norm{\curv_P}/\kappa_{\max,P}$ makes precise the intuition that legal
and mathematical language tolerate less semantic perturbation than literary
language---not because of stricter budgets but because of genuinely different
geometric structures on the semantic bundle.

What began as a synonym randomizer has converged on a gauge-theoretic framework
for controlled semantic transformation.  The synonym replacement machinery is
one implementation detail; the framework accommodates any operator that maps
documents to documents, admits a semantic cost, and can be checked for horizontal
motion under the active profile connection.  As embedding models improve, as
connection learning from trajectory-level supervision matures, and as structured
semantic representations become standard, the architecture is positioned to
scale from the proof-of-concept demonstrated here to a robust, application-grade
perturbation engine applicable to language, cognition, software evolution,
version control, scientific theories, and any other domain whose objects admit
a projection onto a semantic manifold.

\newpage

% ════════════════════════════════════════════════════════════════════════════
% Bibliography
% ════════════════════════════════════════════════════════════════════════════

\begin{thebibliography}{99}

% ── Conceptual Foundations ───────────────────────────────────────────────────

\bibitem{lee2013smooth}
Lee, J.~M. (2013).
\newblock \emph{Introduction to Smooth Manifolds}, 2nd ed.
\newblock Springer.

\bibitem{docarmo1992riemannian}
do~Carmo, M.~P. (1992).
\newblock \emph{Riemannian Geometry}.
\newblock Birkh\"{a}user.

\bibitem{amari2016information}
Amari, S. (2016).
\newblock \emph{Information Geometry and Its Applications}.
\newblock Springer.

\bibitem{liberzon2012optimal}
Liberzon, D. (2012).
\newblock \emph{Calculus of Variations and Optimal Control Theory:
  A Concise Introduction}.
\newblock Princeton University Press.

\bibitem{bertsekas1995dynamic}
Bertsekas, D.~P. (1995).
\newblock \emph{Dynamic Programming and Optimal Control}.
\newblock Athena Scientific.

% ── Diffusion, Colored Noise, and Spectral Structure ─────────────────────────

\bibitem{davidson2026colorednoisediffusionsampling}
Davidson, H., Issachar, N., \& Benaim, S. (2026).
\newblock Colored noise diffusion sampling.
\newblock \emph{arXiv preprint arXiv:2605.30332}.

\bibitem{ho2020ddpm}
Ho, J., Jain, A., \& Abbeel, P. (2020).
\newblock Denoising diffusion probabilistic models.
\newblock \emph{Advances in Neural Information Processing Systems}, 33.

\bibitem{rahaman2019spectralbias}
Rahaman, N., et al. (2019).
\newblock On the spectral bias of neural networks.
\newblock \emph{Proceedings of ICML 2019}.

\bibitem{falck2025fourierspace}
Falck, F., et al. (2025).
\newblock A Fourier space perspective on diffusion models.
\newblock \emph{arXiv preprint arXiv:2505.11278}.

\bibitem{wang2025rotationaltrajectories}
Wang, Y., \& Vastola, J. (2025).
\newblock Geometric structure of diffusion sampling trajectories.
\newblock \emph{arXiv preprint}.

\bibitem{huang2024bluenoise}
Huang, X., et al. (2024).
\newblock Blue noise for diffusion models.
\newblock \emph{SIGGRAPH 2024}.

\bibitem{lipman2022flowmatching}
Lipman, Y., et al. (2022).
\newblock Flow matching for generative modeling.
\newblock \emph{arXiv preprint arXiv:2210.02747}.

\bibitem{albergo2022stochastic}
Albergo, M., \& Vanden-Eijnden, E. (2022).
\newblock Building normalizing flows with stochastic interpolants.
\newblock \emph{arXiv preprint arXiv:2209.15571}.

% ── Semantic Representation and Embeddings ───────────────────────────────────

\bibitem{mikolov2013efficient}
Mikolov, T., Chen, K., Corrado, G., \& Dean, J. (2013).
\newblock Efficient estimation of word representations in vector space.
\newblock \emph{arXiv preprint arXiv:1301.3781}.

\bibitem{devlin2019bert}
Devlin, J., Chang, M.-W., Lee, K., \& Toutanova, K. (2019).
\newblock BERT: Pre-training of deep bidirectional transformers for language
  understanding.
\newblock \emph{Proceedings of NAACL 2019}.

\bibitem{reimers2019sentencebert}
Reimers, N., \& Gurevych, I. (2019).
\newblock Sentence-BERT: Sentence embeddings using Siamese BERT-networks.
\newblock \emph{Proceedings of EMNLP 2019}, 3982--3992.

\bibitem{ethayarajh2019contextual}
Ethayarajh, K. (2019).
\newblock How contextual are contextualized word representations?
\newblock \emph{Proceedings of EMNLP 2019}, 55--65.

\bibitem{larsen2016autoencoding}
Larsen, A.~B.~L., et al. (2016).
\newblock Autoencoding beyond pixels using a learned similarity metric.
\newblock \emph{Proceedings of ICML 2016}.

% ── Language Generation, Paraphrasing, and Augmentation ─────────────────────

\bibitem{miller1995wordnet}
Miller, G.~A. (1995).
\newblock WordNet: A lexical database for English.
\newblock \emph{Communications of the ACM}, 38(11), 39--41.

\bibitem{sennrich2016backtranslation}
Sennrich, R., Haddow, B., \& Birch, A. (2016).
\newblock Improving neural machine translation models with monolingual data.
\newblock \emph{Proceedings of ACL 2016}, 86--96.

\bibitem{raffel2020t5}
Raffel, C., et al. (2020).
\newblock Exploring the limits of transfer learning with a unified
  text-to-text transformer.
\newblock \emph{Journal of Machine Learning Research}, 21(140), 1--67.

\bibitem{brown2020gpt3}
Brown, T., et al. (2020).
\newblock Language models are few-shot learners.
\newblock \emph{Advances in Neural Information Processing Systems}, 33.

\bibitem{hokamp2017lexically}
Hokamp, C., \& Liu, Q. (2017).
\newblock Lexically constrained decoding for sequence generation using grid
  beam search.
\newblock \emph{Proceedings of ACL 2017}, 1535--1546.

\bibitem{wieting2017revisiting}
Wieting, J., \& Gimpel, K. (2017).
\newblock Revisiting recurrent networks for paraphrastic sentence embeddings.
\newblock \emph{Proceedings of ACL 2017}, 2078--2088.

% ── Evaluation ───────────────────────────────────────────────────────────────

\bibitem{zhang2020bertscore}
Zhang, T., et al. (2020).
\newblock BERTScore: Evaluating text generation with BERT.
\newblock \emph{Proceedings of ICLR 2020}.

\bibitem{sellam2020bleurt}
Sellam, T., Das, D., \& Parikh, A. (2020).
\newblock BLEURT: Learning robust metrics for text generation.
\newblock \emph{Proceedings of ACL 2020}.

\end{thebibliography}

\end{document}
